\section*{ПРИЛОЖЕНИЕ А. Развёрнутое сравнение альтернативных структур LPM}

\subsection*{А.1. Зачем нужен отдельный сравнительный раздел}

Одной из типичных слабых сторон прикладных магистерских работ
является ситуация, когда автор сразу показывает выбранное решение и
кратко отмечает, что «другие варианты менее удобны». Для инженерного
отчёта этого иногда достаточно, но для диссертации требуется более
строгая аргументация. Необходимо показать не только то, что выбранный
вариант работает, но и то, что выбор был сделан в пространстве
реальных альтернатив и по понятным критериям.

В настоящей работе сравнение альтернатив особенно важно, потому что
выбор структуры хранения напрямую определяет архитектурную модель.
Если выбрать структуру, идеально подходящую только для lookup, то
можно получить великолепные микропоказатели чтения и провалиться на
update stream. Если выбрать структуру, идеально подходящую только для
bulk replace, то можно сохранить быстрый \texttt{Dump}, но так и не
перейти к потоковой актуализации. Поэтому сравнительный анализ нужен
не как формальность, а как основа всей линии рассуждения.

\subsection*{А.2. Критерии сравнения}

Для рассматриваемой задачи предлагается использовать следующие
критерии:

\begin{enumerate}
\item скорость точечного longest prefix match;
\item стоимость вставки одного префикса;
\item стоимость удаления одного префикса;
\item стоимость полной выгрузки всей таблицы;
\item компактность хранения в памяти;
\item удобство конкурентной обработки;
\item сложность реализации и сопровождения в микросервисе общего
назначения.
\end{enumerate}

Эти критерии выбраны так, чтобы отражать как алгоритмические, так и
системные свойства структуры. Например, высокая скорость lookup
ничего не говорит о том, насколько легко реализовать безопасный
поток подписчиков или дешёвую ресинхронизацию состояния.

\subsection*{А.3. Простые массивы и отсортированные слайсы}

Наиболее простая идея -- хранить маршруты в линейном массиве или
слайсе. В самом примитивном случае поиск требует линейного прохода. В
более развитом варианте можно держать массив отсортированным и
использовать вспомогательные индексы. Такой подход имеет одно
сильное качество: полный \texttt{Dump} чрезвычайно дёшев, поскольку
данные уже лежат в последовательной форме.

Именно к этой семье по operational profile близок legacy backend.
Его преимущества очевидны:

\begin{itemize}
\item очень дешёвый полный снимок;
\item простая модель памяти;
\item предсказуемое внешнее представление.
\end{itemize}

Однако фундаментальный недостаток также очевиден: локальное
обновление редко бывает по-настоящему локальным. Чтобы сохранить
упорядоченность, консистентность и вспомогательные индексы,
приходится перестраивать заметную часть состояния. Для batch-модели
это допустимо, для streaming-модели~-- нет.

\subsection*{А.4. Хеш-таблицы}

Чистая хеш-таблица даёт очень привлекательные свойства для точного
match по ключу. Вставка и чтение в среднем близки к $O(1)$, код
относительно прост, а локальное обновление действительно локально.
Но longest prefix match не является точным match по ключу. Для
поддержки LPM поверх хеш-таблицы нужно либо хранить множество
производных ключей \cite{waldvogel1997scalable}, либо перебирать
длины масок и проверять кандидатов, либо иметь дополнительную
структуру навигации.

Кроме того, чистая hash map плохо подходит для полного обхода в
управляемом упорядоченном виде. Все элементы можно перечислить, но
их естественный порядок не связан с префиксным пространством. Это
особенно неудобно, если \texttt{Dump} должен быть детерминированным
и пригодным для downstream-сравнения.

Таким образом, чистая hash map является хорошим вспомогательным
элементом, но плохим единственным источником истины для данной
задачи.

\subsection*{А.5. Patricia/radix-подходы}

Patricia \cite{morrison1968patricia} и её современные производные
исторически являются очень естественным выбором для prefix-oriented
данных. Их главный плюс -- структурное соответствие задаче longest
prefix match. В дереве легко навигировать по битам или байтам адреса,
а компрессия путей делает структуру экономичнее полного trie.

Сильные стороны Patricia/radix:

\begin{itemize}
\item естественная поддержка префиксной организации;
\item хорошая концептуальная совместимость с LPM;
\item понятная теория и богатая историческая база.
\end{itemize}

Слабые стороны в прикладной постановке могут проявляться в нескольких
местах:

\begin{itemize}
\item не всегда оптимальное использование памяти в реальных in-memory
рантаймах;
\item зависимость производительности от конкретной реализации и
layout;
\item необходимость дополнительной инженерной работы для эффективного
\texttt{Dump} и concurrency management.
\end{itemize}

Если бы задача сводилась только к построению статичного индекса для
prefix lookup, Patricia была бы очень сильным претендентом. Но при
переходе к многоклиентскому сервису с потоковыми обновлениями
возрастают требования к управляемости реализации как системы, а не
только как структуры данных.

\subsection*{А.6. LC-trie}

LC-trie \cite{nilsson1999lctries} полезно рассматривать как
интеллектуальную эволюцию compressed trie, ориентированную на
ускорение поиска. За счёт level compression уменьшается эффективная
глубина навигации, а значит и число условных шагов поиска.

Для данной работы LC-trie особенно ценно как контрпример слишком
узкой инженерной логики. Если бы целью было только уменьшить latency
lookup, именно LC-trie и родственные структуры заслуживали бы
приоритетного внимания. Однако здесь встаёт вопрос о цене локальных
обновлений, о полном дампе и о простоте интеграции в многопоточный
сервис. На этих осях преимущество LC-trie уже не столь очевидно и
существенно зависит от конкретной реализации.

\subsection*{А.7. Controlled prefix expansion}

Подход controlled prefix expansion \cite{srinivasan1999prefixexpansion}
демонстрирует, что иногда можно ускорять LPM ценой контролируемого
увеличения объёма данных. Это особенно полезно в сценариях, где
память относительно дёшева, а latency lookup критична. Для
forwarding-plane устройств и специализированных таблиц это может
быть очень разумным компромиссом.

В прикладном сервисе соответствия \textit{prefix~$\rightarrow$~ASN}
такой подход интересен скорее как теоретический ориентир. Он
показывает важный принцип: не всегда стоит минимизировать размер
структуры любой ценой; иногда правильнее заплатить памятью за
скорость чтения. Однако в потоковой модели этот подход должен быть
оценён ещё и по цене обновлений. Если expansion-схема делает
локальную вставку дорогой, она снова оказывается не в центре
целевого инженерного компромисса.

\subsection*{А.8. Hash-based LPM и binary search on prefix lengths}

Работы Waldvogel и других авторов \cite{waldvogel1997scalable,
gupta1998hardware} показывают, что longest prefix match можно
реализовать не только деревьями, но и комбинацией хеширования,
бинарного поиска по длинам префиксов и вспомогательных структур. Для
среды, где главным врагом являются лишние memory accesses, эти
подходы чрезвычайно убедительны.

Но у них есть тонкое ограничение в рамках данной темы: они
прекрасны для статического или слабо изменяемого forwarding-plane
представления, однако в сервисе, который одновременно должен:

\begin{itemize}
\item принимать поток обновлений;
\item отдавать ordered \texttt{Dump};
\item поддерживать подписки;
\item интегрироваться в обычную серверную кодовую базу,
\end{itemize}

такие структуры становятся сложнее в эксплуатации и анализе. Иными
словами, выигрыш в чистой lookup-метрике может быть перекрыт ростом
сложности состояния как инженерного артефакта.

\subsection*{А.9. Tree Bitmap}

Tree Bitmap \cite{eatherton2004treebitmap} интересен тем, что
пытается совместить достойную скорость lookup, компактность и
поддержку инкрементальных обновлений. Именно это делает его очень
важной альтернативой для обсуждения. По духу он ближе к нашей задаче,
чем чисто hardware-prefix-expansion схемы.

Тем не менее Tree Bitmap в работе не выбран по двум причинам.
Во-первых, его основная литература и инженерная мотивация всё же
сильнее привязаны к маршрутизаторным lookup-сценариям, чем к
общесервисной mutable state модели. Во-вторых, интеграционная
сложность и цена разработки/поддержки в данной кодовой базе для
ART-подобного решения оказались ниже.

Это не означает, что Tree Bitmap хуже по существу. Это означает, что
он менее естественно вписывается в конкретную комбинацию ограничений
текущего проекта.

\subsection*{А.10. Poptrie}

Poptrie \cite{asai2015poptrie} является особенно сильным современным
software baseline. Он показывает, насколько далеко можно продвинуть
lookup-скорость на CPU общего назначения, если агрессивно учитывать
layout данных и возможности процессорных инструкций (population
count). Для работ, где основная цель -- быстрый lookup на полном
маршрутизаторном столе, Poptrie представляет собой мощный ориентир.

Однако у Poptrie есть тот же важный вопрос, что и у других
lookup-first структур: насколько удобно на его основе строить
general-purpose mutable service state? Если основная оптимизация
достигается плотной компрессией и cache-aware layout, то локальное
обновление и concurrent mutation могут становиться более дорогими с
инженерной точки зрения.

Поэтому Poptrie полезен как аргумент для раздела «почему новое
решение не является абсолютным лидером по чтению», но не как
естественный каркас для текущей system architecture.

\subsection*{А.11. ART в сравнительной матрице}

Если собрать сравнительную матрицу в качественном виде, получается
картина, представленная в табл.~\ref{tab:lpm-matrix}.

\begin{table}[ht]
\centering
\caption{Качественное сравнение структур хранения по ключевым критериям}
\label{tab:lpm-matrix}
\small
\begin{tabular}{>{\raggedright\arraybackslash}p{3.0cm}cccccc}
\toprule
Структура         & Lookup & Insert & Delete & Dump  & Memory & Concurrency \\
\midrule
Линейный массив   & низкий & низкий & низкий & \textbf{высокий} & ср.   & низкий \\
Хеш-таблица       & высокий$^{\star}$ & \textbf{высокий} & \textbf{высокий} & ср.   & ср.    & ср.    \\
Patricia/radix    & высокий & ср.    & ср.    & ср.   & ср.    & ср.    \\
LC-trie           & \textbf{высокий} & низкий & низкий & ср.   & ср.    & низкий \\
Tree Bitmap       & \textbf{высокий} & ср.    & ср.    & ср.   & \textbf{высокий} & ср. \\
Poptrie           & \textbf{высокий} & низкий & низкий & ср.   & \textbf{высокий} & низкий \\
\textbf{ART (шард.)} & ср.   & \textbf{высокий} & \textbf{высокий} & ср.   & \textbf{высокий} & \textbf{высокий} \\
\bottomrule
\end{tabular}\\[2pt]
{\footnotesize $^{\star}$ Только для exact match; для LPM требуется
дополнительная схема (например, binary search on prefix lengths).}
\end{table}

Матрица показывает, что ART:

\begin{itemize}
\item не лучший во всех lookup-сценариях;
\item но очень силён как mutable ordered in-memory index;
\item хорошо переносится в сервисную реализацию;
\item допускает шардирование и параллелизм;
\item удобен как основа для локальных updates и full traversal.
\end{itemize}

Именно поэтому ART оказывается не «лучшей LPM-структурой вообще», а
«наиболее подходящим ядром для данного сочетания требований».

\subsection*{А.12. Вывод по сравнительному анализу}

После проведения такого сравнения выбор ART перестаёт выглядеть
интуитивным. Он становится следствием формальной логики:

\begin{enumerate}
\item целевая система должна поддерживать потоковую актуализацию;
\item потоковая актуализация требует дешёвого локального изменения
состояния;
\item дешёвое локальное изменение состояния требует mutable
структуры;
\item mutable структура должна при этом поддерживать prefix-oriented
lookup и ordered \texttt{Dump};
\item ART даёт наиболее практичный компромисс под этот набор
требований.
\end{enumerate}

Следовательно, выбор ART не является случайным. Он вытекает из
постановки задачи и из сравнительного анализа альтернатив.

%==============================================================
